%% ma: L11-B07-0014
%% lop: 11
%% chuong: 2
%% bai: 7
%% dang: 11.7.10
%% loai: DS
%% muc: [TH, TH, NB, NB]
%% dap_an: "SĐĐS"
%% nguon: "de-on-tap-giua-ki-1-toan-11-nam-2026-2027-truong-thpt-ngo-quyen-da-nang.pdf (M09, Đề ôn tập giữa kì 1 Toán 11 THPT Ngô Quyền, Đà Nẵng), Mã đề 0114, Phần II Câu 2"
%% kien_thuc: "Bài 6 (cấp số cộng), Bài 7 (cấp số nhân)"
%% trang_thai: chinh-thuc
%% ghi_chu: "Đã duyệt 29/10/2026."
\begin{ex}
Một nhà hát có $25$ hàng ghế với $16$ ghế ở hàng thứ nhất, $18$ ghế ở hàng thứ hai, $20$ ghế ở hàng thứ ba và cứ tiếp tục theo quy luật đó, tức là hàng sau nhiều hơn hàng liền trước nó $2$ ghế. Khi đó:
\choiceTF{Hàng ghế thứ $10$ có $35$ ghế}{\True Tổng số ghế có trong nhà hát là $1000$ ghế}{\True Số ghế ở mỗi hàng của nhà hát lập thành một cấp số cộng với số hạng đầu $u_1=16$, công sai $d=2$}{Số ghế ở mỗi hàng của nhà hát lập thành một cấp số nhân với số hạng đầu $u_1=16$, công bội $q=2$}
\loigiai{a) Sai: $u_{10}=16+9\cdot2=34$.

b) Đúng: $S_{25}=\dfrac{25}2(32+24\cdot2)=1000$.

c) Đúng.

d) Sai: $\dfrac{18}{16}\ne\dfrac{20}{18}$, không phải cấp số nhân.}
\end{ex}
